% Phase 2 handout -- RISC-V RV32I Datapath Components in Verilog
% EEL 4768 Computer Architecture, Fall 2026
%
% NOTE TO INSTRUCTOR (delete before distributing, or keep as a LaTeX comment):
% This handout is written to match phase_1.tex's structure and boxes. Every
% port list, signal semantic, legality rule and don't-care condition below was
% read directly out of the testbenches and the harness under
% phase_2/source_test/ (config.py, grade_test.py, rtl/alu_tb.v, rtl/decoder_tb.v,
% rtl/imm_tb.v, rtl/rf_bypass_tb.v, rtl/rf_no_bypass_tb.v) and out of the
% Verilog rule checker in phase_2/docker/python-ece552/. Nothing here is
% assumed. Eight items need your sign-off before this goes out:
%
%   1. RUBRIC / SCOPE MISMATCH (Section 2). phase_2/source_test/ grades FIVE
%      testbenches -- 22 checks: 10 ALU ops, 9 decoder groups, imm, rf x2 --
%      but the two Gradescope-facing graders (docker/submit.py and
%      local/grade_project2.py) still only grade FOUR (alu_tb, imm_tb,
%      rf_no_bypass_tb, rf_bypass_tb) at 3/4 = 0.75 each; neither one knows
%      decoder_tb.v exists, and neither ships it. This handout assumes the
%      decoder IS graded and prints 5 x 0.6. Either add the decoder grading
%      block to both graders (assign_even_scores() repoints automatically) or
%      cut Section 8 and the decoder row from the rubric.
%
%   2. FILENAMES. run_test.sh's error text tells students to submit
%      "alu.v, decode.v and rf.v" -- it omits imm.v, and says decode.v while
%      decoder_tb.v instantiates a module named `decoder'. docker/submit.py's
%      presence check is ["alu.v", "imm.v", "rf.v", "project2.txt"]. This
%      handout tells students decoder.v / module decoder. source_test collects
%      by extension so the local run is unaffected either way, but please make
%      run_test.sh's message and submit.py's list agree with Section 1.
%
%   3. ecall VS ebreak. Phase 1 requires every program to terminate with
%      `ecall'. decoder_tb.v puts 32'h00000073 (ecall) in the ILLEGAL group and
%      makes `ebreak' (32'h00100073) the only legal SYSTEM instruction, driving
%      o_halt. As written a phase-1 program cannot run on a phase-2 decoder.
%      Pick one convention before phase 3 connects them.
%
%   4. decoder_tb.v's header comment says JAL/branch operand selects are
%      "unconstrained", but the code has care_op_sel = uses_alu | uses_cmp,
%      and uses_cmp is branches -- so branches ARE checked and must drive
%      o_op1_sel = 0 and o_op2_sel = 0. Section 8 documents the code's actual
%      behavior. Please fix the comment (or the code) so they agree.
%
%   5. project2.txt. Both Gradescope graders nag for a `project2.txt' inherited
%      from the UW-Madison parent repo. Nothing in this course asks for it and
%      this handout does not mention it. Drop the check or define the file.
%
%   6. THE STYLE CHECKER IS NOT IN THE LOCAL LOOP. docker/submit.py runs
%      `ece552 validate -w *.v' fail-closed -- a violation scores 0 before
%      iverilog ever runs -- but source_test/ never runs it. A student can go
%      22/22 locally and score 0 on submission. Section 4 warns them and tells
%      them to run it by hand; consider wiring it into run_test.sh instead.
%      Separately, validators/operator_validator.py's _shift_illegal() has a
%      stray `print(amount.kind)' debug line that leaks onto stdout.
%
%   7. NO SKELETONS EXIST. Unlike phase 1 there is no phase_2/skeletons/. Every
%      port name, width and direction in Sections 5-8 was reverse-engineered
%      from the testbenches' instantiations. If you ship skeletons, diff them
%      against these tables first.
%
%   8. phase_2/submission/ (run_test.sh's default submission directory) is not
%      in the repository. Students have to create it, which Section 9 says.

\documentclass[11pt]{article}

\usepackage[letterpaper,margin=1in]{geometry}
\usepackage{amsmath,amssymb}
\usepackage{booktabs}
\usepackage{array}
\usepackage{multirow}
\usepackage{caption}
\usepackage{xcolor}
\usepackage{enumitem}
\usepackage{tikz}
\usetikzlibrary{arrows.meta,positioning,calc}
\usepackage{listings}
\usepackage[most]{tcolorbox}
\usepackage{fancyhdr}
\usepackage{parskip}
\usepackage[colorlinks=true,linkcolor=blue!50!black,urlcolor=blue!50!black]{hyperref}

\setlist[itemize]{topsep=2pt,itemsep=2pt,parsep=0pt}
\setlist[enumerate]{topsep=2pt,itemsep=2pt,parsep=0pt}

\lstset{
  basicstyle=\ttfamily\small,
  frame=single,
  breaklines=true,
  columns=fullflexible,
  keepspaces=true,
  showstringspaces=false,
  tabsize=4,
  backgroundcolor=\color{black!3}
}

\tcbuselibrary{skins,breakable}

\newtcolorbox{important}{
  colback=red!5!white, colframe=red!55!black, breakable,
  fonttitle=\bfseries, title=Note
}
\newtcolorbox{clarification}{
  colback=blue!5!white, colframe=blue!45!black, breakable,
  fonttitle=\bfseries, title=Clarification
}
\newtcolorbox{pitfall}{
  colback=orange!8!white, colframe=orange!65!black, breakable,
  fonttitle=\bfseries, title=Common Pitfall
}
\newtcolorbox{worked}{
  colback=green!5!white, colframe=green!35!black, breakable,
  fonttitle=\bfseries, title=Verilog Sketch
}
\newtcolorbox{example}{
  colback=green!5!white, colframe=green!35!black, breakable,
  fonttitle=\bfseries, title=Breakdown
}

\setlength{\headheight}{14pt}
\emergencystretch=1.5em

\pagestyle{fancy}
\fancyhf{}
\fancyhead[L]{EEL 4768 -- Computer Architecture}
\fancyhead[R]{Phase 2: RV32I Module Design}
\fancyfoot[C]{\thepage}
\renewcommand{\headrulewidth}{0.4pt}

\title{\vspace{-1.5cm}\textbf{Phase 2}\\[2pt]\Large RISC-V RV32I Module Design in Verilog}
\author{EEL 4768 Computer Architecture \\ Fall 2026}
\date{}

\begin{document}
\maketitle
\thispagestyle{fancy}
\vspace{-1cm}

\section*{Learning Objectives}
By the end of this phase you should be able to:
\begin{itemize}
  \item Write structural and dataflow Verilog for both combinational and sequential logic.
  \item Build the four initial blocks detailed in this phase.
  \item Debug your own hardware designs to fully test them, considering all possible inputs.
\end{itemize}

\begin{important}
You must write your own Verilog test benches to verify your designs. An example test bench will be provided, as well as skeletons for each of the modules.
\\
You also should reference the RV32I reference card from webcourses for the instruction formats, opcodes, and immediate values.
\end{important}

\section{Submission}
\begin{itemize}
  \item \textbf{One submission per group.} You will manually
        submit your work to Webcourses as a zip file. Students in multiple
        groups, or not in a group by the time of submission, will receive
        \textbf{no credit} for the submission.
  \item \textbf{Files to submit}:
\begin{lstlisting}
alu.v
imm.v
rf.v
decoder.v
\end{lstlisting}
        One module per file, and name the module the same as the filename. The testbenches instantiate modules
        named \texttt{alu}, \texttt{imm}, \texttt{rf} and \texttt{decoder} by
        name.
\end{itemize}

\begin{important}
  Skeletons for each verilog file can be found in \texttt{phase\_2/skeletons/}. You must use these skeletons, or the grader will reject your submission.
\end{important}

\section{Grading Rubric}
\label{sec:rubric}
\begin{center}
\begin{tabular}{@{}llr@{}}
\toprule
\textbf{Category} & \textbf{Testbench} & \textbf{Points} \\
\midrule
ALU                  & \texttt{alu\_tb.v}          & 0.75 \\
Immediate Generator  & \texttt{imm\_tb.v}          & 0.75 \\
\multirow[t]{2}{*}{Register File}
                     & \texttt{rf.v (no bypass)} & 0.25 \\
                     & \texttt{rf.v (with bypass)} & 0.5 \\
Instruction Decoder  & \texttt{decoder\_tb.v}      & 0.75 \\
\midrule
\textbf{Total} & & \textbf{3.0} \\
\bottomrule
\end{tabular}
\end{center}

\section{Verilog Coding Rules}
\label{sec:rules}
Your submission is internally checked against a synthesizable subset. This is to insure your code will be synthasizable for future phases. The rules listed below will result in a 0 for the specific file.

\begin{center}
\begin{tabular}{@{}p{5.4cm}p{8.8cm}@{}}
\toprule
\textbf{Rule} & \textbf{What it means for you} \\
\midrule
No \texttt{*}, \texttt{/}, \texttt{\%}, \texttt{**}
  & Arithmetic that does not synthesize cheaply. You do not need any of them
    in this phase. \\
No \texttt{===}, \texttt{!==}, \texttt{==?}, \texttt{!=?}
  & Not synthesizable. Use \texttt{==} and \texttt{!=}. \\
\textbf{Shifts only by a constant}
  & \texttt{x << 3} is fine; \texttt{x << i\_op2[4:0]} is \emph{rejected}. See
    Section~\ref{sec:alu}. \\
No \texttt{if}/\texttt{else} in combinational logic
  & \texttt{if}/\texttt{else} is only legal inside a clocked
    \texttt{always @(posedge \dots)} block. For combinational logic use a
    continuous \texttt{assign} with the \texttt{?:} operator, or a
    \texttt{case} inside \texttt{always @(*)}. \\
No procedural loops
  & \texttt{for}, \texttt{while}, \texttt{repeat} and \texttt{do}/\texttt{while}
    statements are rejected. To build something repetitive, use a
    \texttt{generate} loop, or just write the instances out. \\
No delays or waits
  & \texttt{\#5}, \texttt{wait}, \texttt{disable} and friends are simulation
    constructs. \\
No system tasks
  & \texttt{\$display}, \texttt{\$finish}, \texttt{\$stop}, \texttt{\$system},
    \texttt{\$readmemh}/\texttt{\$readmemb},
    \texttt{\$writememh}/\texttt{\$writememb} and all file I/O
    (\texttt{\$fopen}, \texttt{\$fwrite}, \dots) are not allowed in your final submission. Only
    constant functions such as \texttt{\$clog2} are allowed. \\
No \texttt{`include}
  & Not needed for importing other files. \\
\bottomrule
\end{tabular}
\end{center}

\newpage

\section{Part 1: ALU}
\label{sec:alu}
For the first part, you will be tasked with designing a 32-bit ALU following the RISC-V instruction set. This is a purely combinational ALU. Each ALU operation should be implemented manually. You cannot use +, -, <<, >>, or any other arithmetic operations inplace of the ALU operations. Follow the table below for how to implement each operation.

\begin{center}
  \begin{tabular}{@{}lp{7.8cm}@{}}
    \toprule
    \textbf{\texttt{Operation}} & \textbf{Implementation} \\
    \texttt{Add / Sub} & \texttt{2's compliment Lookahead Carry Adder} \\
    \texttt{SLL, SRL, SRA} & \texttt{Barrel Shifter} \\
    \texttt{SLT / SLTU} & \texttt{Comparator} \\
    \texttt{XOR, OR, AND} & \texttt{Bitwise Logic} \\
    \bottomrule
  \end{tabular}
\end{center}

Below lists the input and output signals, and what they represent.

\begin{center}
\begin{tabular}{@{}llp{8.4cm}@{}}
\toprule
\textbf{Port} & \textbf{Dir/Width} & \textbf{Meaning} \\
\midrule
\texttt{i\_opsel}    & in \ [2:0]  & Selects the operation (table below). \\
\texttt{i\_sub}      & in          & With \texttt{i\_opsel = 3'b000}, subtract instead of add. Ignored for every other \texttt{i\_opsel}. \\
\texttt{i\_unsigned} & in          & Comparisons are unsigned. Only used for branch comparisons. Unsigned arithmetic does not use this.\\
\texttt{i\_arith}    & in          & With \texttt{i\_opsel = 3'b101}, shift right \emph{arithmetically} (replicate the sign bit) instead of logically. \\
\texttt{i\_op1}      & in \ [31:0] & First operand. The value shifted, for shifts. \\
\texttt{i\_op2}      & in \ [31:0] & Second operand. \textbf{For shifts, only \texttt{i\_op2[4:0]} is the shift amount}; the upper 27 bits are ignored. \\
\texttt{o\_result}   & out [31:0]  & The operation's result. \\
\texttt{o\_eq}       & out         & \texttt{i\_op1 == i\_op2}. \\
\texttt{o\_slt}      & out         & \texttt{i\_op1 < i\_op2}, signed or unsigned per \texttt{i\_unsigned}. \\
\bottomrule
\end{tabular}
\end{center}

\begin{center}
\begin{tabular}{@{}llp{7.8cm}@{}}
\toprule
\textbf{\texttt{i\_opsel}} & \textbf{Operation} & \textbf{\texttt{o\_result}} \\
\midrule
\texttt{3'b000} & \textsc{add} / \textsc{sub} & \texttt{i\_op1 + i\_op2}, or \texttt{i\_op1 - i\_op2} when \texttt{i\_sub}. \\
\texttt{3'b001} & \textsc{sll}                & \texttt{i\_op1} shifted left by \texttt{i\_op2[4:0]}. \\
\texttt{3'b010} & \textsc{slt} & \texttt{i\_op1 < i\_op2}, signed. \\
\texttt{3'b011} & \textsc{sltu} & \texttt{i\_op1 < i\_op2}, unsigned. \\
\texttt{3'b100} & \textsc{xor}                & \texttt{i\_op1 \^{} i\_op2}. \\
\texttt{3'b101} & \textsc{srl} / \textsc{sra} & \texttt{i\_op1} shifted right by \texttt{i\_op2[4:0]}, logically or (when \texttt{i\_arith}) arithmetically. \\
\texttt{3'b110} & \textsc{or}                 & \texttt{i\_op1 | i\_op2}. \\
\texttt{3'b111} & \textsc{and}                & \texttt{i\_op1 \& i\_op2}. \\
\bottomrule
\end{tabular}
\end{center}

\begin{clarification}
\textbf{\texttt{o\_eq} and \texttt{o\_slt} are always checked.} Even if they are not directly used, it is important to always insure they output.
\end{clarification}

\newpage

\section{Part 2: Immediate Generator}
\label{sec:imm}
Part 2 implements an immediate generator. This is acombinational block that extracts and sign-extends the immediate value from an 32-bit instruction word.

\begin{center}
\begin{tabular}{@{}llp{8.6cm}@{}}
\toprule
\textbf{Port} & \textbf{Dir/Width} & \textbf{Meaning} \\
\midrule
\texttt{i\_inst}      & in \ [31:0] & The full 32-bit instruction word. \\
\texttt{i\_format}    & in \ [5:0]  & One-hot Selector for the instruction format. \\
\texttt{o\_immediate} & out [31:0]  & The 32-bit sign-extended immediate. \\
\bottomrule
\end{tabular}
\end{center}

\begin{center}
\begin{tabular}{@{}lll@{}}
\toprule
\textbf{\texttt{i\_format}} & \textbf{Bit} & \textbf{Format} \\
\midrule
\texttt{6'b000001} & [0] & R \\
\texttt{6'b000010} & [1] & I \\
\texttt{6'b000100} & [2] & S \\
\texttt{6'b001000} & [3] & B \\
\texttt{6'b010000} & [4] & U \\
\texttt{6'b100000} & [5] & J \\
\bottomrule
\end{tabular}
\captionof{table}{The six instruction formats and the one-hot
\texttt{i\_format} code that selects each.}
\label{tab:imm}
\end{center}

\newpage

\section{Part 3: Register File}
\label{sec:rf}
This part implements a 32-bit register file, or the registers of your processor. Reads are combinational, while writes are sequential.

\begin{center}
\begin{tabular}{@{}llp{8.2cm}@{}}
\toprule
\textbf{Port} & \textbf{Dir/Width} & \textbf{Meaning} \\
\midrule
\texttt{i\_clk}       & in          & Clock. Everything sequential happens on its rising edge. \\
\texttt{i\_rst}       & in          & \textbf{Synchronous}, active high. Clears all 32 registers to zero. \\
\texttt{i\_rs1\_raddr}& in \ [4:0]  & Read port 1 address. \\
\texttt{o\_rs1\_rdata}& out [31:0]  & Read port 1 data. \\
\texttt{i\_rs2\_raddr}& in \ [4:0]  & Read port 2 address. \\
\texttt{o\_rs2\_rdata}& out [31:0]  & Read port 2 data.\\
\texttt{i\_rd\_wen}   & in          & Write enable. \\
\texttt{i\_rd\_waddr} & in \ [4:0]  & Write address. \\
\texttt{i\_rd\_wdata} & in \ [31:0] & Write data. \\
\bottomrule
\end{tabular}
\end{center}

The module takes \textbf{one parameter}, the bypass enable, and it must be the
\emph{first} parameter --- the two testbenches instantiate it positionally, as
\texttt{rf \#(0)} and \texttt{rf \#(1)}:

\begin{lstlisting}[language=Verilog]
module rf #(parameter BYPASS_EN = 0) (
    input  wire        i_clk,
    /* ... */
);
\end{lstlisting}

The bypass enable allows the register file to enact read-before-write behavior, so that a read port can see the new value of a register on the same clock edge it is written. You have to handle both cases for full points.



\subsection{Behavior}
\begin{itemize}
  \item \textbf{\texttt{x0} is hardwired to zero.} A write to address
        \texttt{5'b00000} must be discarded, and a read of address
        \texttt{5'b00000} must return \texttt{32'b0} forever after. Both
        testbenches write \texttt{32'hDEADBEEF} to \texttt{x0} as their very
        first check.
  \item \textbf{Writes are synchronous}, taking effect on the rising edge of
        \texttt{i\_clk} when \texttt{i\_rd\_wen} is high.
  \item \textbf{Reads are combinational.} \texttt{o\_rs1\_rdata} tracks
        \texttt{i\_rs1\_raddr} with no clock edge in between.
  \item \textbf{Reset is synchronous and active high}, applied inside the same
        clocked block as the write.
  \item \textbf{Bypass}, when \texttt{BYPASS\_EN} is nonzero: if a write is in
        flight (\texttt{i\_rd\_wen} is high) to the same register a read port
        is addressing, and that register is not \texttt{x0}, that read port
        returns \texttt{i\_rd\_wdata} \emph{immediately}, without waiting for
        the clock edge. When \texttt{BYPASS\_EN} is zero the read port returns
        the stored value, and the write is only visible after the edge.
\end{itemize}

\newpage

\section{Part 4: Instruction Decoder}
\label{sec:decoder}
The final part has you constructing the instruction decoder. This is a
combinational block that takes one 32-bit instruction word and produces every control signal the datapath needs.

\begin{important}
  The decoder instantiates uses generator you wrote in Section~\ref{sec:imm}. You must instantiate the imm module within your decoder.
\end{important}

\subsection{Instruction fields}
Everything the decoder does is a function of these fields:

\begin{center}
\begin{tabular}{@{}llp{8.6cm}@{}}
\toprule
\textbf{Field} & \textbf{Bits} & \textbf{Use} \\
\midrule
\texttt{opcode} & \texttt{i\_inst[6:0]}   & Selects the instruction class. \\
\texttt{rd}     & \texttt{i\_inst[11:7]}  & Destination register specifier. \\
\texttt{funct3} & \texttt{i\_inst[14:12]} & Sub-operation within a class. \\
\texttt{rs1}    & \texttt{i\_inst[19:15]} & First source register specifier. \\
\texttt{rs2}    & \texttt{i\_inst[24:20]} & Second source register specifier. \\
\texttt{funct7} & \texttt{i\_inst[31:25]} & Further sub-operation (add vs.\ sub, srl vs.\ sra). \\
\bottomrule
\end{tabular}
\end{center}


\subsection{Outputs}
\label{sec:decoder-outputs}

\noindent\begin{minipage}{\linewidth}
\textbf{Legality and halt.}.

\smallskip

\begin{center}
\begin{tabular}{@{}llp{12.0cm}@{}}
\toprule
\textbf{Port} & \textbf{Width} & \textbf{Meaning} \\
\midrule
\texttt{o\_legal} & 1 & The word is a legal RV32I encoding. \\
\texttt{o\_halt}  & 1 & The instruction is \texttt{ebreak} (\texttt{32'h00100073}). \\
\bottomrule
\end{tabular}
\end{center}
\end{minipage}
\par\medskip

\begin{important}
\textbf{\texttt{ecall} is \emph{not} legal in this decoder.} \texttt{ebreak}
(\texttt{32'h00100073}) is the only \textsc{system} instruction, and it is what
drives \texttt{o\_halt}. \texttt{ecall} (\texttt{32'h00000073}) is tested in the
\textsc{illegal} group and must produce \texttt{o\_legal = 0} and
\texttt{o\_halt = 0}. This differs from Phase 1, where programs terminate with
\texttt{ecall}.
\end{important}

\noindent\begin{minipage}{\linewidth}
\textbf{Register specifiers and immediate.}

\smallskip

\begin{center}
\begin{tabular}{@{}llp{12.0cm}@{}}
\toprule
\textbf{Port} & \textbf{Width} & \textbf{Meaning} \\
\midrule
\texttt{o\_rs1}       & [4:0]  & \texttt{i\_inst[19:15]}, for instructions that read \texttt{rs1}. \\
\texttt{o\_rs2}       & [4:0]  & \texttt{i\_inst[24:20]}, for instructions that read \texttt{rs2} (\textsc{op}, \textsc{store}, \textsc{branch}). \\
\texttt{o\_rd}        & [4:0]  & \texttt{i\_inst[11:7]} for instructions that write a register, and \textbf{\texttt{5'd0} for every instruction that does not} --- stores, branches, and any illegal word. \\
\texttt{o\_immediate} & [31:0] & The immediate for this instruction's format (Table~\ref{tab:imm} lists the format codes). Instantiate your \texttt{imm} module here. \\
\bottomrule
\end{tabular}
\end{center}
\end{minipage}
\par\medskip

\noindent\begin{minipage}{\linewidth}
\textbf{ALU control.}

\smallskip

\begin{center}
\begin{tabular}{@{}llp{12.0cm}@{}}
\toprule
\textbf{Port} & \textbf{Width} & \textbf{Meaning} \\
\midrule
\texttt{o\_op1\_sel}     & 1     & ALU operand 1 source: \texttt{0} = \texttt{rs1}, \texttt{1} = PC. Only \texttt{auipc} selects the PC. \\
\texttt{o\_op2\_sel}     & 1     & ALU operand 2 source: \texttt{0} = \texttt{rs2}, \texttt{1} = immediate. Set for \textsc{op-imm}, \textsc{load}, \textsc{store}, \texttt{auipc} and \texttt{jalr}. \\
\texttt{o\_alu\_opsel}   & [2:0] & \texttt{alu}'s \texttt{i\_opsel}: \texttt{funct3} for \textsc{op} and \textsc{op-imm}, and \texttt{3'b000} (add) for everything else that uses the ALU. \\
\texttt{o\_alu\_sub}     & 1     & \texttt{alu}'s \texttt{i\_sub}. Set \emph{only} for \texttt{sub} (\textsc{op}, \texttt{funct3} \texttt{000}, \texttt{funct7} \texttt{0100000}); clear for every address calculation. \\
\texttt{o\_alu\_unsigned}& 1     & \texttt{alu}'s \texttt{i\_unsigned}. Set for \texttt{sltu}/\texttt{sltiu}, and equal to \texttt{funct3[1]} for branches (so \texttt{bltu}/\texttt{bgeu} set it). \\
\texttt{o\_alu\_arith}   & 1     & \texttt{alu}'s \texttt{i\_arith}, \texttt{funct7 == 7'b0100000}, for \texttt{sra}/\texttt{srai}. \\
\bottomrule
\end{tabular}
\end{center}
\end{minipage}
\par\medskip

\noindent\begin{minipage}{\linewidth}
\textbf{Branch and jump control.}

\smallskip

\begin{center}
\begin{tabular}{@{}llp{12.0cm}@{}}
\toprule
\textbf{Port} & \textbf{Width} & \textbf{Meaning} \\
\midrule
\texttt{o\_branch}          & 1 & This is a (legal) conditional branch. \\
\texttt{o\_jump}            & 1 & This is \texttt{jal} or \texttt{jalr}. \\
\texttt{o\_branch\_equal}   & 1 & The branch tests equality (\texttt{beq}, \texttt{bne}) rather than ordering: \texttt{\~{}funct3[2]}. \\
\texttt{o\_branch\_unsigned}& 1 & The ordering comparison is unsigned (\texttt{bltu}, \texttt{bgeu}): \texttt{funct3[1]}. \\
\texttt{o\_branch\_invert}  & 1 & Take the branch on the \emph{negation} of the test (\texttt{bne}, \texttt{bge}, \texttt{bgeu}): \texttt{funct3[0]}. \\
\texttt{o\_pc\_sel}         & 1 & Target base: \texttt{0} = PC (branches, \texttt{jal}), \texttt{1} = \texttt{rs1} (\texttt{jalr}). \\
\bottomrule
\end{tabular}
\end{center}
\end{minipage}
\par\medskip

\noindent\begin{minipage}{\linewidth}
\textbf{Data-memory control.}

\smallskip

\begin{center}
\begin{tabular}{@{}llp{13.0cm}@{}}
\toprule
\textbf{Port} & \textbf{Width} & \textbf{Meaning} \\
\midrule
\texttt{o\_dmem\_ren}   & 1     & This is a load. \\
\texttt{o\_dmem\_wen}   & 1     & This is a store. \\
\texttt{o\_dmem\_align} & [1:0] & Alignment mask of the access: \texttt{2'b00} byte, \texttt{2'b01} half-word, \texttt{2'b11} word. \\
\texttt{o\_dmem\_memb}  & 1     & Byte access (\texttt{lb}, \texttt{lbu}, \texttt{sb}). \\
\texttt{o\_dmem\_memh}  & 1     & Half-word access (\texttt{lh}, \texttt{lhu}, \texttt{sh}). \\
\texttt{o\_dmem\_memw}  & 1     & Word access (\texttt{lw}, \texttt{sw}). \\
\texttt{o\_dmem\_memu}  & 1     & The loaded value is \emph{zero}-extended rather than sign-extended (\texttt{lbu}, \texttt{lhu}). \\
\bottomrule
\end{tabular}
\end{center}
\end{minipage}
\par\medskip

\textbf{Writeback select.} \texttt{o\_rd\_sel [3:0]} is \textbf{one-hot} and
chooses what gets written to \texttt{rd}:

\begin{center}
\begin{tabular}{@{}llll@{}}
\toprule
\textbf{Bit} & \textbf{Value} & \textbf{Source} & \textbf{Instructions} \\
\midrule
{}[0] & \texttt{4'b0001} & ALU result & \textsc{op}, \textsc{op-imm}, \texttt{auipc} \\
{}[1] & \texttt{4'b0010} & immediate  & \texttt{lui} \\
{}[2] & \texttt{4'b0100} & PC + 4     & \texttt{jal}, \texttt{jalr} \\
{}[3] & \texttt{4'b1000} & memory     & \textsc{load} \\
\bottomrule
\end{tabular}
\end{center}



\newpage

\section{Testing Your Work}
\label{sec:testing}
\textbf{There is no autograder in this repository.} Verifying that
\texttt{alu.v}, \texttt{imm.v}, \texttt{rf.v} and \texttt{decoder.v} behave
correctly is part of the assignment.

\subsection{The example}
\texttt{phase\_2/example/} holds two files:

\begin{itemize}
\item \texttt{opmux.v} --- a small combinational module with four inputs
  (\texttt{i\_a}, \texttt{i\_b}, \texttt{i\_sel}, \texttt{i\_en}) and two
  outputs (\texttt{o\_result}, \texttt{o\_zero}). A two-bit select mux chooses
  between \texttt{+}, \texttt{-}, \texttt{<<} and \texttt{>>}.
\item \texttt{opmux\_tb.v} --- a self-checking testbench for it. It is commented as a walkthrough to help you build your own test benches.
\end{itemize}

Run it with \texttt{./phase\_2/run\_example.sh}, which prints the two commands
it uses so you can run them on your own modules:

\begin{lstlisting}[language={}]
$ iverilog -s opmux_tb -o build/sim example/opmux_tb.v example/opmux.v
$ ./build/sim

========== opmux testbench ==========
--- add ---
[PASS] add: 7 + 9
[PASS] add: 0 + 0 sets zero
...
212 passed, 0 failed
ALL TESTS PASSED
\end{lstlisting}

\texttt{-s} names the top module to elaborate, \texttt{-o} names the simulator
to write, and every source file the design needs is listed after them.

\begin{important}
\textbf{An expected value copied out of your design proves only that the
design agrees with itself.} Make sure to fully test all edge cases of your own design.
\end{important}

\end{document}